%%
% Copyright (c) 2017 - 2024, Pascal Wagler;
% Copyright (c) 2014 - 2024, John MacFarlane
%
% All rights reserved.
%
% Redistribution and use in source and binary forms, with or without
% modification, are permitted provided that the following conditions
% are met:
%
% - Redistributions of source code must retain the above copyright
% notice, this list of conditions and the following disclaimer.
%
% - Redistributions in binary form must reproduce the above copyright
% notice, this list of conditions and the following disclaimer in the
% documentation and/or other materials provided with the distribution.
%
% - Neither the name of John MacFarlane nor the names of other
% contributors may be used to endorse or promote products derived
% from this software without specific prior written permission.
%
% THIS SOFTWARE IS PROVIDED BY THE COPYRIGHT HOLDERS AND CONTRIBUTORS
% "AS IS" AND ANY EXPRESS OR IMPLIED WARRANTIES, INCLUDING, BUT NOT
% LIMITED TO, THE IMPLIED WARRANTIES OF MERCHANTABILITY AND FITNESS
% FOR A PARTICULAR PURPOSE ARE DISCLAIMED. IN NO EVENT SHALL THE
% COPYRIGHT OWNER OR CONTRIBUTORS BE LIABLE FOR ANY DIRECT, INDIRECT,
% INCIDENTAL, SPECIAL, EXEMPLARY, OR CONSEQUENTIAL DAMAGES (INCLUDING,
% BUT NOT LIMITED TO, PROCUREMENT OF SUBSTITUTE GOODS OR SERVICES;
% LOSS OF USE, DATA, OR PROFITS; OR BUSINESS INTERRUPTION) HOWEVER
% CAUSED AND ON ANY THEORY OF LIABILITY, WHETHER IN CONTRACT, STRICT
% LIABILITY, OR TORT (INCLUDING NEGLIGENCE OR OTHERWISE) ARISING IN
% ANY WAY OUT OF THE USE OF THIS SOFTWARE, EVEN IF ADVISED OF THE
% POSSIBILITY OF SUCH DAMAGE.
%%

%%
% This is the Eisvogel pandoc LaTeX template.
%
% For usage information and examples visit the official GitHub page:
% https://github.com/Wandmalfarbe/pandoc-latex-template
%%

% Options for packages loaded elsewhere
\PassOptionsToPackage{unicode}{hyperref}
\PassOptionsToPackage{hyphens}{url}
\PassOptionsToPackage{dvipsnames,svgnames,x11names,table}{xcolor}
%
\documentclass[
  paper=a4,
  oneside  ,captions=tableheading
]{scrartcl}
\usepackage{amsmath,amssymb}
% Use setspace anyway because we change the default line spacing.
% The spacing is changed early to affect the titlepage and the TOC.
\usepackage{setspace}
\setstretch{1.2}
\usepackage{iftex}
\ifPDFTeX
  \usepackage[T1]{fontenc}
  \usepackage[utf8]{inputenc}
  \usepackage{textcomp} % provide euro and other symbols
\else % if luatex or xetex
  \usepackage{unicode-math} % this also loads fontspec
  \defaultfontfeatures{Scale=MatchLowercase}
  \defaultfontfeatures[\rmfamily]{Ligatures=TeX,Scale=1}
\fi
\usepackage{lmodern}
\ifPDFTeX\else
  % xetex/luatex font selection
    \setmainfont[]{Inter}
    \setmonofont[]{FiraCode}
\fi
% Use upquote if available, for straight quotes in verbatim environments
\IfFileExists{upquote.sty}{\usepackage{upquote}}{}
\IfFileExists{microtype.sty}{% use microtype if available
  \usepackage[]{microtype}
  \UseMicrotypeSet[protrusion]{basicmath} % disable protrusion for tt fonts
}{}
\makeatletter
\@ifundefined{KOMAClassName}{% if non-KOMA class
  \IfFileExists{parskip.sty}{%
    \usepackage{parskip}
  }{% else
    \setlength{\parindent}{0pt}
    \setlength{\parskip}{6pt plus 2pt minus 1pt}}
}{% if KOMA class
  \KOMAoptions{parskip=half}}
\makeatother
\usepackage{xcolor}
\definecolor{default-linkcolor}{HTML}{A50000}
\definecolor{default-filecolor}{HTML}{A50000}
\definecolor{default-citecolor}{HTML}{4077C0}
\definecolor{default-urlcolor}{HTML}{4077C0}
\usepackage[margin=2.5cm,includehead=true,includefoot=true,centering,]{geometry}
\newcounter{none}
\usepackage{longtable,booktabs,array}
\usepackage{calc} % for calculating minipage widths
\usepackage{caption}
% Correct order of tables after \paragraph or \subparagraph
\usepackage{etoolbox}
\makeatletter
\patchcmd\longtable{\par}{\if@noskipsec\mbox{}\fi\par}{}{}
\makeatother
% Allow footnotes in longtable head/foot
\IfFileExists{footnotehyper.sty}{\usepackage{footnotehyper}}{\usepackage{footnote}}
\makesavenoteenv{longtable}
% add backlinks to footnote references, cf. https://tex.stackexchange.com/questions/302266/make-footnote-clickable-both-ways
\usepackage{footnotebackref}
\usepackage{graphicx}
\makeatletter
\newsavebox\pandoc@box
\newcommand*\pandocbounded[1]{% scales image to fit in text height/width
  \sbox\pandoc@box{#1}%
  \Gscale@div\@tempa{\textheight}{\dimexpr\ht\pandoc@box+\dp\pandoc@box\relax}%
  \Gscale@div\@tempb{\linewidth}{\wd\pandoc@box}%
  \ifdim\@tempb\p@<\@tempa\p@\let\@tempa\@tempb\fi% select the smaller of both
  \ifdim\@tempa\p@<\p@\scalebox{\@tempa}{\usebox\pandoc@box}%
  \else\usebox{\pandoc@box}%
  \fi%
}
% Set default figure placement to htbp
% Make use of float-package and set default placement for figures to H.
% The option H means 'PUT IT HERE' (as  opposed to the standard h option which means 'You may put it here if you like').
\usepackage{float}
\floatplacement{figure}{H}
\makeatother
\setlength{\emergencystretch}{3em} % prevent overfull lines
\providecommand{\tightlist}{%
  \setlength{\itemsep}{0pt}\setlength{\parskip}{0pt}}
\setcounter{secnumdepth}{5}
% Make \paragraph and \subparagraph free-standing
\makeatletter
\ifx\paragraph\undefined\else
  \let\oldparagraph\paragraph
  \renewcommand{\paragraph}{
    \@ifstar
      \xxxParagraphStar
      \xxxParagraphNoStar
  }
  \newcommand{\xxxParagraphStar}[1]{\oldparagraph*{#1}\mbox{}}
  \newcommand{\xxxParagraphNoStar}[1]{\oldparagraph{#1}\mbox{}}
\fi
\ifx\subparagraph\undefined\else
  \let\oldsubparagraph\subparagraph
  \renewcommand{\subparagraph}{
    \@ifstar
      \xxxSubParagraphStar
      \xxxSubParagraphNoStar
  }
  \newcommand{\xxxSubParagraphStar}[1]{\oldsubparagraph*{#1}\mbox{}}
  \newcommand{\xxxSubParagraphNoStar}[1]{\oldsubparagraph{#1}\mbox{}}
\fi
\makeatother
\ifLuaTeX
\usepackage[bidi=basic]{babel}
\else
\usepackage[bidi=default]{babel}
\fi
\babelprovide[main,import]{catalan}
\ifPDFTeX
\else
\babelfont{rm}[]{Inter}
\fi
% get rid of language-specific shorthands (see #6817):
\let\LanguageShortHands\languageshorthands
\def\languageshorthands#1{}
\usepackage{fontspec}

\setmainfont[
  Path=_extensions/eisvogel/fonts/Inter/static/,
  Extension=.ttf,
  UprightFont=*-Regular,
  BoldFont=*-Bold,
  ItalicFont=*-Italic,
  BoldItalicFont=*-BoldItalic
]{Inter}

\setsansfont[
  Path=_extensions/eisvogel/fonts/Inter/static/,
  Extension=.ttf,
  UprightFont=*-Regular,
  BoldFont=*-Bold,
  ItalicFont=*-Italic,
  BoldItalicFont=*-BoldItalic
]{Inter}

\setmonofont[
  Path=_extensions/eisvogel/fonts/Fira_Code/static/,
  Extension=.ttf,
  UprightFont=*-Regular,
  BoldFont=*-Bold
]{FiraCode}
\makeatletter
\@ifpackageloaded{caption}{}{\usepackage{caption}}
\AtBeginDocument{%
\ifdefined\contentsname
  \renewcommand*\contentsname{Taula de continguts}
\else
  \newcommand\contentsname{Taula de continguts}
\fi
\ifdefined\listfigurename
  \renewcommand*\listfigurename{Llista de figures}
\else
  \newcommand\listfigurename{Llista de figures}
\fi
\ifdefined\listtablename
  \renewcommand*\listtablename{Llista de taules}
\else
  \newcommand\listtablename{Llista de taules}
\fi
\ifdefined\figurename
  \renewcommand*\figurename{Figura}
\else
  \newcommand\figurename{Figura}
\fi
\ifdefined\tablename
  \renewcommand*\tablename{Taula}
\else
  \newcommand\tablename{Taula}
\fi
}
\@ifpackageloaded{float}{}{\usepackage{float}}
\floatstyle{ruled}
\@ifundefined{c@chapter}{\newfloat{codelisting}{h}{lop}}{\newfloat{codelisting}{h}{lop}[chapter]}
\floatname{codelisting}{Llistat}
\newcommand*\listoflistings{\listof{codelisting}{Llista de llistats}}
\makeatother
\makeatletter
\makeatother
\makeatletter
\@ifpackageloaded{caption}{}{\usepackage{caption}}
\@ifpackageloaded{subcaption}{}{\usepackage{subcaption}}
\makeatother
\usepackage{bookmark}
\IfFileExists{xurl.sty}{\usepackage{xurl}}{} % add URL line breaks if available
\urlstyle{same}
\hypersetup{
  pdftitle={Projecte de millora de l'eficiencia del fred a Euromed},
  pdfauthor={Arnau Coronado},
  pdflang={ca},
  colorlinks=true,
  linkcolor={blue},
  filecolor={default-filecolor},
  citecolor={default-citecolor},
  urlcolor={default-urlcolor},
  breaklinks=true,
  pdfcreator={LaTeX via pandoc with the Eisvogel template}}
\title{Projecte de millora de l'eficiencia del fred a Euromed}
\author{Arnau Coronado}
\date{2026-11-17}



%%
%% added
%%


%
% for the background color of the title page
%
\usepackage{pagecolor}
\usepackage{afterpage}
\usepackage[margin=2.5cm,includehead=true,includefoot=true,centering]{geometry}

%
% break urls
%
\PassOptionsToPackage{hyphens}{url}

%
% When using babel or polyglossia with biblatex, loading csquotes is recommended
% to ensure that quoted texts are typeset according to the rules of your main language.
%
\usepackage{csquotes}

%
% captions
%
\definecolor{caption-color}{HTML}{777777}
\usepackage[font={stretch=1.2}, textfont={color=caption-color}, position=top, skip=4mm, labelfont=bf, singlelinecheck=false, justification=raggedright]{caption}
\setcapindent{0em}

%
% blockquote
%
\definecolor{blockquote-border}{RGB}{221,221,221}
\definecolor{blockquote-text}{RGB}{119,119,119}
\usepackage{mdframed}
\newmdenv[rightline=false,bottomline=false,topline=false,linewidth=3pt,linecolor=blockquote-border,skipabove=\parskip]{customblockquote}
\renewenvironment{quote}{\begin{customblockquote}\list{}{\rightmargin=0em\leftmargin=0em}%
\item\relax\color{blockquote-text}\ignorespaces}{\unskip\unskip\endlist\end{customblockquote}}

%
% Source Sans Pro as the default font family
% Source Code Pro for monospace text
%
% 'default' option sets the default
% font family to Source Sans Pro, not \sfdefault.
%
\ifnum 0\ifxetex 1\fi\ifluatex 1\fi=0 % if pdftex
    % \usepackage[default]{sourcesanspro}
  \usepackage{sourcecodepro}
  \else % if not pdftex
    \fi

%
% heading color
%
\definecolor{heading-color}{RGB}{40,40,40}
\addtokomafont{section}{\color{heading-color}}
% When using the classes report, scrreprt, book,
% scrbook or memoir, uncomment the following line.
%\addtokomafont{chapter}{\color{heading-color}}

%
% variables for title, author and date
%
\usepackage{titling}
\title{Projecte de millora de l'eficiencia del fred a Euromed}
\author{Arnau Coronado}
\date{2026-11-17}

%
% tables
%

\definecolor{table-row-color}{HTML}{F5F5F5}
\definecolor{table-rule-color}{HTML}{14612C}

%\arrayrulecolor{black!40}
\arrayrulecolor{table-rule-color}     % color of \toprule, \midrule, \bottomrule
\setlength\heavyrulewidth{0.3ex}      % thickness of \toprule, \bottomrule
\renewcommand{\arraystretch}{1.3}     % spacing (padding)
\AtBeginEnvironment{longtable}{\fontsize{6.5pt}{7.5pt}\selectfont}


%
% remove paragraph indentation
%
\setlength{\parindent}{0pt}
\setlength{\parskip}{6pt plus 2pt minus 1pt}
\setlength{\emergencystretch}{3em}  % prevent overfull lines

%
%
% Listings
%
%


%
% header and footer
%
\usepackage[headsepline,footsepline]{scrlayer-scrpage}

\newpairofpagestyles{eisvogel-header-footer}{
  \clearpairofpagestyles
  \ihead*{Euromed S.A.}
  \chead*{}
  \ohead*{2026-11-17}
  \ifoot*{Arnau Coronado}
  \cfoot*{}
  \ofoot*{\thepage}
  \addtokomafont{pageheadfoot}{\upshape}
}
\pagestyle{eisvogel-header-footer}



%%
%% end added
%%

\begin{document}

%%
%% begin titlepage
%%
\begin{titlepage}
\newgeometry{left=6cm}
\newcommand{\colorRule}[3][black]{\textcolor[HTML]{#1}{\rule{#2}{#3}}}
\begin{flushleft}
\noindent
\\[-1em]
\color[HTML]{000000}
\makebox[0pt][l]{\colorRule[14612c]{1.3\textwidth}{4pt}}
\par
\noindent

{
  \setstretch{1.4}
  \vfill
  \noindent {\huge \textbf{\textsf{Projecte de millora de l'eficiencia
del fred a Euromed}}}
    \vskip 2em
  \noindent {\Large \textsf{Arnau Coronado}}
  \vfill
}


\textsf{2026-11-17}
\end{flushleft}
\end{titlepage}
\restoregeometry
\pagenumbering{arabic}

%%
%% end titlepage
%%

% \maketitle


\renewcommand*\contentsname{Taula de continguts}
{
\hypersetup{linkcolor=}
\setcounter{tocdepth}{3}
\tableofcontents
\newpage
}
\section{Introducció}\label{introducciuxf3}

Euromed, S.A. és una empresa industrial amb seu a Mollet del Vallès
(Barcelona), especialitzada en el desenvolupament i producció
d'extractes botànics per a la indústria farmacèutica, nutracèutica i
alimentària. L'empresa forma part del grup Euromed {[}@Euromed2024{]},
amb presència internacional i orientació a mercats europeus, americans i
asiàtics.

La companyia disposa d'instal·lacions industrials avançades amb
processos d'extracció, purificació i assecat {[}@Euromed2024{]}, que
impliquen l'ús intensiu de fluids tèrmics, aigua de procés i sistemes
auxiliars. Aquestes instal·lacions requereixen un control rigorós dels
paràmetres operatius per garantir tant la qualitat del producte com
l'eficiència energètica.

Des del punt de vista operatiu, la planta integra múltiples circuits de
fluids (aigua calenta, aigua freda, vapor, etc.) que intervenen
directament en els processos productius. El control dels cabals en
aquests circuits és un factor crític per a l'optimització del consum
energètic i la detecció de possibles desviacions operatives.

\section{Antecedents}\label{antecedents}

Aquest projecte actual es la suma (o desembocadura) de dos projectes
exectuats des de la primavera de 2025. A continuació es resumiran els
resultats obtinguts en aquests estudis.

\subsection{Projecte ``Remodelació circuits de fred'' de
DeDietrich}\label{projecte-remodelaciuxf3-circuits-de-fred-de-dedietrich}

El projecte de De Dietrich va tenir com a objectiu obtenir una visió
conjunta dels circuits de fred de la fàbrica i identificar oportunitats
per millorar-ne el funcionament. La documentació recollida a la carpeta
\texttt{DDIB} constitueix la base de l'estudi i inclou la revisió i
actualització dels diagrames de canonades i instrumentació (P\&ID) dels
circuits d'aigua de torres i d'aigua refrigerada.

L'estudi també va inventariar els equips que consumeixen fred i va
reunir, per a cada equip, les dades disponibles de volum, cabal i
temperatures. Això permet relacionar els diferents consumidors amb els
circuits que els abasteixen, entendre com es distribueix l'aigua i
disposar d'una referència comuna per analitzar possibles desequilibris o
limitacions de servei.

Partint d'aquesta representació de la instal·lació i de les dades dels
equips, De Dietrich va plantejar diverses alternatives d'actuació.
Aquestes propostes s'han de llegir en el context i amb els condicionants
descrits als documents de \texttt{DDIB}; no totes necessàriament es van
executar. El treball va servir, sobretot, per establir una diagnosi i
orientar les decisions posteriors sobre els circuits. El projecte
SuperSònic complementa aquest estudi amb mesures de cabal reals, que
permeten contrastar les hipòtesis i valorar les millores proposades.

\subsection{Projecte SuperSònic}\label{projecte-supersuxf2nic}

El projecte SuperSònic acompanya la inversió en un cabalímetre
ultrasònic no invasiu FLUXUS G608, de FLEXIM. Té com a objectiu mesurar
el cabal en diversos punts dels circuits de fred de la fàbrica i obtenir
dades reals del seu funcionament. El treball parteix de l'estudi de De
Dietrich, que aporta el coneixement dels circuits i dels seus
consumidors; les mesures de SuperSònic, al seu torn, permeten contrastar
les hipòtesis plantejades i valorar les millores proposades.

Les dades es consulten a la base de dades del projecte i s'organitzen
per data i punt de mesura. Si hi ha més d'una lectura per al mateix punt
i dia, el tractament actual conserva el valor màxim. A partir de les
lectures també es calculen alguns cabals derivats, emprant relacions
entre punts de mesura i una estimació del nombre de bombes en servei. En
l'anàlisi dels resultats, per tant, cal diferenciar les mesures directes
del cabalímetre dels valors calculats i tenir en compte els supòsits
utilitzats en aquests càlculs.

\subsubsection{Resultats destacats de
SuperSònic}\label{resultats-destacats-de-supersuxf2nic}

El resum següent es genera a partir de les lectures disponibles a la
base de dades. La primera taula recull punts mesurats directament; la
segona, indicadors calculats a partir de les relacions i supòsits
descrits anteriorment. Els valors són cabals en les unitats registrades
pel sistema.

\begin{verbatim}
Període analitzat: 2026-09-07 – 2026-11-04
Dies amb dades: 16
Punts de mesura directa amb les medianes més elevades:
\end{verbatim}

{\def\LTcaptype{none} % do not increment counter
\begin{longtable}[]{@{}llll@{}}
\toprule\noalign{}
& Dies amb lectura & Mediana & Màxim \\
measurement\_point & & & \\
\midrule\noalign{}
\endhead
\bottomrule\noalign{}
\endlastfoot
05.06 & 14.0 & 667.72 & 828.46 \\
06.07 & 14.0 & 650.66 & 799.92 \\
8108 & 8.0 & 279.64 & 294.44 \\
8107 & 8.0 & 237.88 & 250.41 \\
8104 & 11.0 & 216.57 & 235.06 \\
\end{longtable}
}

\begin{verbatim}
Indicadors calculats (no són lectures directes del cabalímetre):
\end{verbatim}

{\def\LTcaptype{none} % do not increment counter
\begin{longtable}[]{@{}lllll@{}}
\toprule\noalign{}
& Dies amb valor & Mediana & Màxim & Mínim \\
measurement\_point & & & & \\
\midrule\noalign{}
\endhead
\bottomrule\noalign{}
\endlastfoot
04.05 & 14.0 & 759.82 & 959.36 & 637.86 \\
01.04 & 12.0 & 802.07 & 982.35 & 742.73 \\
Nautas & 12.0 & 46.82 & 114.88 & -2.12 \\
06.22 & 14.0 & 20.84 & 73.16 & -15.37 \\
06.10 & 14.0 & 650.66 & 799.92 & 574.91 \\
5701 & 13.0 & 52.26 & 178.65 & 18.71 \\
3800.01 & 15.0 & 102.21 & 135.84 & 80.35 \\
31001 & 13.0 & 106.66 & 127.42 & 49.33 \\
3901 & 12.0 & 10.95 & 16.76 & -7.88 \\
11.19 & 14.0 & 174.21 & 203.16 & 153.34 \\
sil-pel & 12.0 & 61.52 & 73.84 & 52.38 \\
11.15 & 9.0 & 27.13 & 29.05 & 20.81 \\
\end{longtable}
}

\section{Actuacions}\label{actuacions}

\subsection{Incendi edifici de la
E-700}\label{incendi-edifici-de-la-e-700}

El calendari de les actuacions ha estat marcat per l'incendi que es va
produir la matinada de l'11 d'abril de 2026 a la fàbrica. L'incendi va
cremar les tres plantes de l'edifici que contenia la planta E-700 i les
columnes R-150 i R-450. El circuit de fred que anava a aquesta zona de
la fabrica es va anular; reduint així el consum de fred global i
augmentant el cabal per la resta de consumidors.

Algunes de les actuacions proposades per DeDietrich que se centraven en
aquesta zona de la fabrica (que tenia un dèficit de pressió ja que els
condensadors estaven situats a molta alçada) van quedar desestimades i
ara la prioritat és la reconstrucció de les plantes malmeses i incloure
les millores a la fase de disseny.

\subsection{Vàlvules E-800, E900 i
E-1000}\label{vuxe0lvules-e-800-e900-i-e-1000}

El dia 21 d'octubre de 2026 es va posar en funcionament les valvules de
la planta E-800, per avaluar-ne el funcionament i, posteriorment,
activar també les de la E900/E1000.

Les vàlvules tallen el pas d'aigua de refrigeració als condensadors
durant les fases de treball que no requereixen fred. S'espera que,
mentre la vàlvula d'un equip estigui tancada, el cabal que hi deixa de
circular es redistribueixi per la xarxa i millori el cabal disponible
per a la resta d'equips.

\begin{itemize}
\tightlist
\item
  Fases evaporador

  \begin{itemize}
  \tightlist
  \item
    Posar en buit
  \item
    Posar a pressió atmosfèrica
  \item
    Evaporació a pressió atmosfèrica
  \item
    Evaporació en buit
  \item
    Concentració en buit
  \item
    Buidat evaporador
  \end{itemize}
\end{itemize}

{\def\LTcaptype{none} % do not increment counter
\begin{longtable}[]{@{}ll@{}}
\toprule\noalign{}
Fases de treball & Valvula oberta/tancada \\
\midrule\noalign{}
\endhead
\bottomrule\noalign{}
\endlastfoot
Posar en buit & Oberta \\
Posar a pressió atmosfèrica & Oberta \\
Evaporació a pressió atmosfèrica & Oberta \\
Evaporació en buit & Oberta \\
Concentració en buit & Oberta \\
Buidat evaporador & Tancada \\
\end{longtable}
}

\begin{itemize}
\tightlist
\item
  Fases extractor

  \begin{itemize}
  \tightlist
  \item
    Inertització amb nitrogen
  \item
    Posar al buit
  \item
    Posar a pressió atmosfèrica
  \item
    Càrrega de sòlids
  \item
    Càrrega de dissolvent
  \item
    Extracció
  \item
    Filtració sota pressió
  \item
    Assecat a pressió atmosfèrica
  \item
    Assecat al buit
  \item
    Descarrega residus
  \end{itemize}
\end{itemize}

{\def\LTcaptype{none} % do not increment counter
\begin{longtable}[]{@{}ll@{}}
\toprule\noalign{}
Fases de treball & Valvula oberta/tancada \\
\midrule\noalign{}
\endhead
\bottomrule\noalign{}
\endlastfoot
Inertització amb nitrogen & Tancada \\
Posar al buit & Oberta \\
Posar a pressió atmosfèrica & Oberta \\
Càrrega de sòlids & Tancada \\
Càrrega de dissolvent & Tancada \\
Extracció & Oberta \\
Filtració sota pressió & Tancada \\
Assecat a pressió atmosfèrica & Oberta \\
Assecat al buit & Oberta \\
Descarrega residus & Tancada \\
\end{longtable}
}

\subsection{Variadors de freqüencia}\label{variadors-de-frequxfcencia}

L'altra actuació que s'ha executat és la instal·lació de variadors de
freqüència a les bombes de l'aigua freda i els ventiladors de les
torres. Aquesta actuació permet ajustar el cabal de les bombes a la
demanda real del sistema i la potència de les torres a la temperatura de
l'aigua, millorant així l'eficiència energètica i reduint el consum
d'electricitat.

\section{Resultats}\label{resultats}

\subsection{Avaluacio de les valvules E-800, E900 i
E-1000}\label{avaluacio-de-les-valvules-e-800-e900-i-e-1000}

Per avaluar les vàlvules de la E-800 s'ha mesurat el cabal al punt 3806,
a l'entrada del condensador de l'evaporador (E-806), amb el cabalímetre
FLUXUS G608. El registre va del 30.10.2026 a les 14:00 fins al
31.10.2026 a les 17:05 (27,1 h, una lectura per minut) i s'ha comparat
amb la fase en curs de l'evaporador i de l'extractor registrada al
SCADA. L'extractor és al costat de l'evaporador, té el seu propi
condensador i la seva pròpia vàlvula, i està connectat a la mateixa
xarxa. A les figures, la fase de cada equip es mostra com a franja de
color darrere la corba de cabal.

\begin{figure}

\centering{

\pandocbounded{\includegraphics[keepaspectratio]{figures/combined/flow_3806_evapora.png}}

}

\caption{\label{fig-flow-3806-evapora}Cabal al punt 3806 (condensador
E-806) i fase en curs de l'evaporador de la E-800.}

\end{figure}%

\begin{figure}

\centering{

\pandocbounded{\includegraphics[keepaspectratio]{figures/combined/flow_3806_extract.png}}

}

\caption{\label{fig-flow-3806-extract}Cabal al punt 3806 (condensador
E-806) i fase en curs de l'extractor de la E-800.}

\end{figure}%

El cabal només adopta tres nivells: 0, 80,8 ± 2,4 m³/h i 101,7 ± 2,7
m³/h (mitjana ± desviació estàndard). La
Taula~\ref{tbl-cabal-estat-valvules} mostra quin nivell es dona en cada
combinació d'estats de les dues vàlvules. L'estat de cada vàlvula s'ha
deduït de la fase en curs segons les taules de l'apartat d'actuacions.
La fase 7 de l'evaporador, en què el cabal és sempre nul, s'ha comptat
com a vàlvula tancada.

\begin{longtable}[]{@{}
  >{\raggedright\arraybackslash}p{(\linewidth - 16\tabcolsep) * \real{0.0882}}
  >{\raggedright\arraybackslash}p{(\linewidth - 16\tabcolsep) * \real{0.0882}}
  >{\raggedleft\arraybackslash}p{(\linewidth - 16\tabcolsep) * \real{0.1176}}
  >{\raggedleft\arraybackslash}p{(\linewidth - 16\tabcolsep) * \real{0.1176}}
  >{\raggedleft\arraybackslash}p{(\linewidth - 16\tabcolsep) * \real{0.1176}}
  >{\raggedleft\arraybackslash}p{(\linewidth - 16\tabcolsep) * \real{0.1176}}
  >{\raggedleft\arraybackslash}p{(\linewidth - 16\tabcolsep) * \real{0.1176}}
  >{\raggedleft\arraybackslash}p{(\linewidth - 16\tabcolsep) * \real{0.1176}}
  >{\raggedleft\arraybackslash}p{(\linewidth - 16\tabcolsep) * \real{0.1176}}@{}}
\caption{Cabal al punt 3806 segons l'estat de les vàlvules de
l'evaporador i de l'extractor (30--31.10.2026). Els percentatges de
nivell de cabal són sobre el temps de cada
fila.}\label{tbl-cabal-estat-valvules}\tabularnewline
\toprule\noalign{}
\begin{minipage}[b]{\linewidth}\raggedright
Vàlvula evaporador
\end{minipage} & \begin{minipage}[b]{\linewidth}\raggedright
Vàlvula extractor
\end{minipage} & \begin{minipage}[b]{\linewidth}\raggedleft
Temps (h)
\end{minipage} & \begin{minipage}[b]{\linewidth}\raggedleft
\% del període
\end{minipage} & \begin{minipage}[b]{\linewidth}\raggedleft
0 m³/h
\end{minipage} & \begin{minipage}[b]{\linewidth}\raggedleft
≈ 80 m³/h
\end{minipage} & \begin{minipage}[b]{\linewidth}\raggedleft
≈ 100 m³/h
\end{minipage} & \begin{minipage}[b]{\linewidth}\raggedleft
Cabal mitjà (m³/h)
\end{minipage} & \begin{minipage}[b]{\linewidth}\raggedleft
Volum (m³)
\end{minipage} \\
\midrule\noalign{}
\endfirsthead
\toprule\noalign{}
\begin{minipage}[b]{\linewidth}\raggedright
Vàlvula evaporador
\end{minipage} & \begin{minipage}[b]{\linewidth}\raggedright
Vàlvula extractor
\end{minipage} & \begin{minipage}[b]{\linewidth}\raggedleft
Temps (h)
\end{minipage} & \begin{minipage}[b]{\linewidth}\raggedleft
\% del període
\end{minipage} & \begin{minipage}[b]{\linewidth}\raggedleft
0 m³/h
\end{minipage} & \begin{minipage}[b]{\linewidth}\raggedleft
≈ 80 m³/h
\end{minipage} & \begin{minipage}[b]{\linewidth}\raggedleft
≈ 100 m³/h
\end{minipage} & \begin{minipage}[b]{\linewidth}\raggedleft
Cabal mitjà (m³/h)
\end{minipage} & \begin{minipage}[b]{\linewidth}\raggedleft
Volum (m³)
\end{minipage} \\
\midrule\noalign{}
\endhead
\bottomrule\noalign{}
\endlastfoot
Oberta & Oberta & 7,1 & 26,2 \% & 0,2 \% & 98,4 \% & 1,4 \% & 81,1 &
575 \\
Oberta & Tancada & 9,3 & 34,3 \% & 12,6 \% & 17,4 \% & 70,0 \% & 85,4 &
792 \\
Tancada & Oberta & 2,0 & 7,5 \% & 80,3 \% & 19,7 \% & 0,0 \% & 15,6 &
32 \\
Tancada & Tancada & 8,7 & 32,0 \% & 96,0 \% & 0,0 \% & 4,0 \% & 4,3 &
37 \\
\textbf{Total} & & \textbf{27,1} & \textbf{100 \%} & \textbf{41,1 \%} &
\textbf{33,2 \%} & \textbf{25,6 \%} & \textbf{53,0} & \textbf{1.436} \\
\end{longtable}

De la taula se'n desprèn el següent:

\begin{itemize}
\tightlist
\item
  \textbf{Vàlvula de l'evaporador tancada.} El cabal és nul la major
  part del temps (el 96 \% quan l'extractor també està tancat). Gairebé
  tot el temps restant correspon als minuts que segueixen cada
  tancament, a causa del retard descrit més avall. En total, el
  condensador de l'evaporador va estar sense cabal un 41 \% del període
  (11,2 h), mentre l'extractor continuava treballant. En les 10,7 h amb
  la vàlvula de l'evaporador tancada, el condensador només va rebre 69
  m³, el 5 \% del volum total.
\item
  \textbf{Les dues vàlvules obertes.} El cabal és estable a 80 m³/h
  (98,4 \% del temps). Aquest és el cabal de referència del condensador
  de l'evaporador.
\item
  \textbf{Vàlvula de l'evaporador oberta i vàlvula de l'extractor
  tancada.} El cabal puja a 100 m³/h (70 \% del temps), un augment d'uns
  21 m³/h (+26 \%). El 17 \% del temps a 80 m³/h correspon sobretot al
  retard de tancament de la vàlvula de l'extractor.
\end{itemize}

Aquest augment no vol dir que tot el cabal que deixa de passar pel
condensador de l'extractor vagi al de l'evaporador. En tancar un ramal
canvia la corba resistent de la xarxa, i el cabal alliberat es reparteix
entre tots els consumidors que continuen oberts. És d'esperar que el més
beneficiat sigui el més proper, en aquest cas l'evaporador, però els 21
m³/h en són només una part (vegeu la comparativa amb el condensador de
l'extractor més avall). Quan les dues vàlvules estan tancades, tot el
cabal es redistribueix cap a la resta de la fàbrica.

La Taula~\ref{tbl-temps-resposta} mostra el temps que triga el cabal a
respondre després de cada maniobra de vàlvula.

\begin{longtable}[]{@{}
  >{\raggedright\arraybackslash}p{(\linewidth - 6\tabcolsep) * \real{0.2143}}
  >{\raggedleft\arraybackslash}p{(\linewidth - 6\tabcolsep) * \real{0.2857}}
  >{\raggedleft\arraybackslash}p{(\linewidth - 6\tabcolsep) * \real{0.2857}}
  >{\raggedright\arraybackslash}p{(\linewidth - 6\tabcolsep) * \real{0.2143}}@{}}
\caption{Temps de resposta del cabal al punt 3806 després de cada
maniobra de vàlvula (resolució ±1
min).}\label{tbl-temps-resposta}\tabularnewline
\toprule\noalign{}
\begin{minipage}[b]{\linewidth}\raggedright
Maniobra
\end{minipage} & \begin{minipage}[b]{\linewidth}\raggedleft
Esdeveniments
\end{minipage} & \begin{minipage}[b]{\linewidth}\raggedleft
Temps de resposta
\end{minipage} & \begin{minipage}[b]{\linewidth}\raggedright
Nivell resultant
\end{minipage} \\
\midrule\noalign{}
\endfirsthead
\toprule\noalign{}
\begin{minipage}[b]{\linewidth}\raggedright
Maniobra
\end{minipage} & \begin{minipage}[b]{\linewidth}\raggedleft
Esdeveniments
\end{minipage} & \begin{minipage}[b]{\linewidth}\raggedleft
Temps de resposta
\end{minipage} & \begin{minipage}[b]{\linewidth}\raggedright
Nivell resultant
\end{minipage} \\
\midrule\noalign{}
\endhead
\bottomrule\noalign{}
\endlastfoot
Obertura de la vàlvula de l'evaporador & 3 & 1,0 -- 1,5 min & 80 m³/h \\
Tancament de la vàlvula de l'evaporador & 3 & 11,1 -- 12,0 min & 0
m³/h \\
Obertura de la vàlvula de l'extractor (evaporador en servei) & 7 & 0,8
-- 1,9 min & 80 m³/h \\
Tancament de la vàlvula de l'extractor (evaporador en servei) & 7 & 10,9
-- 11,4 min & 100 m³/h \\
\end{longtable}

La resposta és molt reproduïble: les vàlvules obren en uns 1--2 minuts i
tanquen sistemàticament uns 11 minuts després del canvi de fase. Això
apunta a una temporització programada al tancament o a un actuador lent.
Aquest comportament evita maniobres curtes, però allarga la circulació
d'aigua després de cada fase que no requereix fred, i convé verificar-lo
a la configuració de l'automatisme. Només hi ha dues anomalies puntuals.
El 31.10 a les 02:30 el cabal passa a ser nul en plena fase 4 de
l'evaporador, sense cap canvi de fase. El 31.10, entre les 06:48 i les
06:58, hi ha un cabal de 100 m³/h amb l'evaporador sense demanda.

El 04.11.2026, de 16:44 a 19:40 (2,9 h, una lectura per segon), es va
mesurar també el cabal al punt 3801, a l'entrada del condensador de
l'extractor (Figura~\ref{fig-flow-3801}). El cabal torna a adoptar tres
nivells discrets: 0 m³/h (0,2 h), 29,8 ± 0,6 m³/h (2,1 h) i 61,0 ± 1,1
m³/h (0,6 h). L'exportació del SCADA disponible acaba el 03.11, de
manera que per a aquest registre no es coneix la fase de cada equip. Per
analogia amb el punt 3806, s'interpreta que el nivell de 0 m³/h
correspon a la vàlvula de l'extractor tancada, el de 30 m³/h a les dues
vàlvules obertes i el de 60 m³/h a la vàlvula de l'extractor oberta amb
la de l'evaporador tancada.

\begin{figure}

\centering{

\pandocbounded{\includegraphics[keepaspectratio]{figures/supersonic/3801_88.9mm_2026-11-04.png}}

}

\caption{\label{fig-flow-3801}Cabal al punt 3801 (condensador de
l'extractor de la E-800), 04.11.2026.}

\end{figure}%

\begin{longtable}[]{@{}
  >{\raggedright\arraybackslash}p{(\linewidth - 10\tabcolsep) * \real{0.1364}}
  >{\raggedright\arraybackslash}p{(\linewidth - 10\tabcolsep) * \real{0.1364}}
  >{\raggedleft\arraybackslash}p{(\linewidth - 10\tabcolsep) * \real{0.1818}}
  >{\raggedleft\arraybackslash}p{(\linewidth - 10\tabcolsep) * \real{0.1818}}
  >{\raggedleft\arraybackslash}p{(\linewidth - 10\tabcolsep) * \real{0.1818}}
  >{\raggedleft\arraybackslash}p{(\linewidth - 10\tabcolsep) * \real{0.1818}}@{}}
\caption{Cabal (m³/h) als condensadors de l'evaporador i de l'extractor
de la E-800 segons l'estat de la seva vàlvula i de la vàlvula de l'equip
homòleg. Les mesures dels dos punts són de dies
diferents.}\label{tbl-comparativa-condensadors}\tabularnewline
\toprule\noalign{}
\begin{minipage}[b]{\linewidth}\raggedright
Condensador
\end{minipage} & \begin{minipage}[b]{\linewidth}\raggedright
Punt
\end{minipage} & \begin{minipage}[b]{\linewidth}\raggedleft
Vàlvula pròpia tancada
\end{minipage} & \begin{minipage}[b]{\linewidth}\raggedleft
Vàlvula pròpia oberta, homòloga oberta
\end{minipage} & \begin{minipage}[b]{\linewidth}\raggedleft
Vàlvula pròpia oberta, homòloga tancada
\end{minipage} & \begin{minipage}[b]{\linewidth}\raggedleft
Increment
\end{minipage} \\
\midrule\noalign{}
\endfirsthead
\toprule\noalign{}
\begin{minipage}[b]{\linewidth}\raggedright
Condensador
\end{minipage} & \begin{minipage}[b]{\linewidth}\raggedright
Punt
\end{minipage} & \begin{minipage}[b]{\linewidth}\raggedleft
Vàlvula pròpia tancada
\end{minipage} & \begin{minipage}[b]{\linewidth}\raggedleft
Vàlvula pròpia oberta, homòloga oberta
\end{minipage} & \begin{minipage}[b]{\linewidth}\raggedleft
Vàlvula pròpia oberta, homòloga tancada
\end{minipage} & \begin{minipage}[b]{\linewidth}\raggedleft
Increment
\end{minipage} \\
\midrule\noalign{}
\endhead
\bottomrule\noalign{}
\endlastfoot
Evaporador (E-806) & 3806 & 0 & 80,8 & 101,7 & +20,9 (+26 \%) \\
Extractor & 3801 & 0 & 29,8 & 61,0 & +31,2 (+105 \%) \\
\end{longtable}

La Taula~\ref{tbl-comparativa-condensadors} confirma el mateix
comportament als dos condensadors. Cada vàlvula talla completament el
pas d'aigua quan el seu equip no necessita fred. Quan la vàlvula
homòloga tanca, el cabal augmenta clarament, i en termes relatius
l'augment és molt més gran a l'extractor, que duplica el seu cabal. Els
valors també permeten estimar quina part del cabal alliberat aprofita
l'equip veí:

\begin{itemize}
\tightlist
\item
  Quan tanca l'extractor deixen de circular-hi uns 30 m³/h, i
  l'evaporador en guanya 21 (≈ 70 \%).
\item
  Quan tanca l'evaporador deixen de circular-hi uns 81 m³/h, i
  l'extractor en guanya 31 (≈ 38 \%).
\end{itemize}

La resta es reparteix per la xarxa. Aquesta estimació és aproximada,
perquè els dos punts es van mesurar en dies diferents i l'estat de la
resta de la xarxa (bombes en servei i altres consumidors) podia ser
diferent.

En conclusió, les vàlvules de la E-800 funcionen d'acord amb la lògica
de fases definida. Cada condensador només rep aigua quan el procés la
necessita, i el cabal que no circula per la E-800 queda disponible per a
la resta de la xarxa. Per confirmar la interpretació del punt 3801 i
saber on va exactament el cabal alliberat, caldria mesurar
simultàniament els dos condensadors i el col·lector principal, amb
l'exportació del SCADA del mateix període. L'avaluació es repetirà quan
s'activin les vàlvules de la E-900 i la E-1000.

\section{Conclusions}\label{conclusions}

\end{document}
